\chapter{Machine-Learning Guidance for Water Quality}

Machine learning in this program is a research tool, not a shortcut around scientific judgment. Models are useful when they help evaluate patterns in water-quality measurements, prioritize follow-up questions, or compare feature sets. Models are not acceptable when they obscure data problems, inflate weak evidence, or produce claims that the sampling design cannot support.

\section{Beginner Script Ladder}

Students should start with the four-script ladder before modifying any model. Each script is intentionally small enough to inspect, and each one writes outputs that can be opened in RStudio.

Run these commands from the \file{training_wicced_water_quality/} folder:

\begin{rcode}
Rscript scripts/04_ml_01_build_feature_table.R
Rscript scripts/04_ml_02_explore_features.R
Rscript scripts/04_ml_03_train_models.R
Rscript scripts/04_ml_04_diagnostics_and_model_card.R
\end{rcode}

\begin{longtable}{p{0.24\textwidth}p{0.34\textwidth}p{0.32\textwidth}}
\toprule
\textbf{Step} & \textbf{What it teaches} & \textbf{Key outputs} \\
\midrule
Build feature table & Targets, features, exclusions, leakage notes, one-hot context variables, and time-based train/test split. & Feature table, feature dictionary, split plan, training rows, testing rows. \\
Explore features & Summary statistics, class balance, simple correlations, and plots before modeling. & Feature summary, correlation table, class balance, exploration plots, student questions. \\
Train models & Mean-only regression baseline, linear regression, majority-class baseline, logistic classifier, metrics, predictions, and coefficients. & Regression metrics, classification metrics, predictions, coefficients, model-comparison summary. \\
Diagnose and document & Prediction plots, residuals, confusion-matrix thinking, model-card writing, and interpretation limits. & Diagnostic plots, diagnostic summary, model card, interpretation guide. \\
\bottomrule
\end{longtable}

\begin{programbox}{Beginner Rule}
The first ML goal is not to build the most accurate model. The first goal is to build a model that can be inspected, compared to a baseline, explained without overclaiming, and rerun by another student.
\end{programbox}

\section{Guiding Resources}

The \file{ml_resources/} folder contains beginner support files:
\begin{itemize}
  \item \file{README.md}: the beginner path and required interpretation rule.
  \item \file{ml_glossary.csv}: plain-language definitions for target, feature, leakage, baseline, residual, model card, and common metrics.
  \item \file{ml_workflow_checklist.csv}: a checklist from question to handoff.
  \item \file{ml_troubleshooting.md}: common beginner problems and how to reason through them.
  \item \file{model_card_field_guide.md}: what to write in each model-card section.
\end{itemize}

\section{What Students Should Open First}

After the four scripts run, students should open these files in order:
\begin{enumerate}
  \item \file{results/ml/ml_feature_dictionary.csv}. Confirm which variables are targets, predictors, excluded variables, and metadata.
  \item \file{results/ml/ml_train_test_split_plan.csv}. Explain what the split tests.
  \item \file{results/ml/ml_exploration_questions.md}. Answer each question in the research log.
  \item \file{results/ml/ml_regression_metrics.csv}. Compare the linear model to the mean-only baseline.
  \item \file{results/ml/ml_classification_metrics.csv}. Compare the logistic model to the majority-class baseline.
  \item \file{results/ml/ml_predictions.csv}. Find rows where the model missed.
  \item \file{results/ml/ml_model_card.md}. Revise the model card for the student's project.
\end{enumerate}

\section{Beginner Feature Policy}

The teaching scripts deliberately exclude \code{conductivity_ms_cm} from the beginner salinity and high-salinity models because conductivity is closely related to salinity. That exclusion teaches leakage caution. Conductivity can be added later only as a documented sensitivity model when the research question says it is available and scientifically appropriate.

\section{Advanced Extension After Beginner Mastery}

The advanced extension is required only after the beginner ladder has been fully understood. It should not be used as a shortcut to make a model look more impressive. Before starting the advanced scripts, students must be able to explain the target, the included features, excluded variables, leakage risks, train/test split, baseline comparison, prediction errors, and the beginner model card.

Run these commands from the \file{training_wicced_water_quality/} folder only after the four beginner ML scripts have been completed:

\begin{rcode}
Rscript scripts/05_ml_01_advanced_validation.R
Rscript scripts/05_ml_02_permutation_importance_and_sensitivity.R
Rscript scripts/05_ml_03_advanced_model_report.R
\end{rcode}

\begin{longtable}{p{0.25\textwidth}p{0.35\textwidth}p{0.30\textwidth}}
\toprule
\textbf{Advanced step} & \textbf{What it adds} & \textbf{Key outputs} \\
\midrule
Blocked validation & Compares the beginner feature set and sensitivity feature sets under a time holdout and leave-one-site-out folds. & Validation folds, split comparison, validation notes. \\
Sensitivity and importance & Tests conductivity as a flagged sensitivity feature, skips overfit feature sets, and estimates held-out permutation importance. & Feature-set metrics, permutation importance, sensitivity warnings, sensitivity plots. \\
Advanced review & Converts advanced outputs into a decision gate and written interpretation. & Advanced model review, decision gate, final recommendations. \\
\bottomrule
\end{longtable}

\begin{cautionbox}{Advanced ML Gate}
An advanced model is not better because it is more complex. It is better only if it improves held-out performance, survives leakage review, remains interpretable, and matches the scientific question. A sensitivity result that improves performance can still be inappropriate for the main analysis.
\end{cautionbox}

The advanced extension writes outputs to the advanced ML results folder. Students should open the model-review markdown file first, then the decision-gate CSV, then the validation and permutation-importance tables. Any gate marked as requiring review must be resolved before advanced ML outputs appear in a final report.

The \file{ml_resources/advanced_ml_extension_guide.md} file gives the student-facing sequence and interpretation rules for this extension.

\section{Acceptable ML Question Types}

\begin{longtable}{p{0.24\textwidth}p{0.34\textwidth}p{0.32\textwidth}}
\toprule
\textbf{Question type} & \textbf{Example} & \textbf{Good first model} \\
\midrule
Regression & Predict salinity, dissolved oxygen, nitrate, or chlorophyll from available environmental features. & Mean-only baseline, linear model, regularized model, or tree model with careful validation. \\
Classification & Classify samples as high-salinity versus lower-salinity using a pre-defined threshold. & Majority-class baseline, logistic regression, decision tree, or random forest if data volume supports it. \\
Anomaly screening & Flag observations that are unusual relative to the student's dataset. & Rule-based flags first; simple distance or residual screening only after QA/QC. \\
Feature-set comparison & Compare water-quality variables alone versus water-quality plus source-backed chemical traits. & Same split and same target across feature sets. \\
Exploratory grouping & Cluster samples or compounds to identify patterns for follow-up. & PCA or clustering as exploratory analysis with no causal claim. \\
\bottomrule
\end{longtable}

\section{Leakage Control}

Leakage occurs when information reaches the model that would not be available at prediction time or that directly encodes the target. Leakage makes performance look better than it is.

Common water-quality leakage examples:
\begin{itemize}
  \item Predicting salinity while including conductivity if the purpose is to predict salinity without a salinity-related sensor. Conductivity may be valid in some designs, but the relationship must be explicit.
  \item Including site ID when the train and test sets contain the same sites and the target is mostly site-specific.
  \item Randomly splitting repeated measurements from the same site/date when the goal is generalization to new dates or sites.
  \item Including post-treatment, post-event, or manually interpreted variables that would not be available during deployment.
  \item Including chemical traits that were selected after inspecting the target response.
\end{itemize}

\section{Split Designs}

\begin{longtable}{p{0.22\textwidth}p{0.34\textwidth}p{0.34\textwidth}}
\toprule
\textbf{Split design} & \textbf{Use when} & \textbf{Caution} \\
\midrule
Random split & Rows are independent and the question is narrow. & Often too optimistic for repeated site/date data. \\
Time-based split & The project asks whether earlier data predict later observations. & Seasonal shifts can dominate performance. \\
Site-holdout split & The project asks whether a model generalizes to new sites. & Requires enough sites and careful interpretation. \\
Group split & Samples are nested by site, date, batch, instrument run, or treatment. & Groups must be defined before modeling. \\
Cross-validation & Data are limited but model comparison is needed. & Folds must respect groups if repeated measures exist. \\
\bottomrule
\end{longtable}

\section{Metrics}

Regression models should report at least one scale-dependent metric and one visual diagnostic:
\begin{itemize}
  \item RMSE: penalizes larger errors.
  \item MAE: average absolute error in the target unit.
  \item Test-set observed-versus-predicted plot.
  \item Residual plot, especially across site/date or predicted value.
\end{itemize}

Classification models should report:
\begin{itemize}
  \item Confusion matrix.
  \item Accuracy only with class balance reported.
  \item Sensitivity and specificity when false negatives and false positives matter.
  \item Threshold used to convert probabilities into classes.
  \item If probabilities are used, a probability distribution or calibration check.
\end{itemize}

\section{Feature Engineering Rules}

Feature engineering should be boring, documented, and defensible. Good features are created from measurements or source-backed records that would be available at prediction time. Transformations should be simple unless there is a strong reason.

Rules:
\begin{itemize}
  \item Freeze the analysis row set before comparing models.
  \item Keep raw variables recoverable.
  \item Record units before scaling or transformation.
  \item Avoid transformations that depend on the whole dataset if the model would later be deployed.
  \item Treat missingness as information only when the missingness mechanism is plausible and documented.
  \item Use source-backed chemical traits only after an evidence audit.
\end{itemize}

\section{Model Cards}

Every model must have a model card. A model card is a short technical record that states:
\begin{itemize}
  \item Project question.
  \item Prediction target.
  \item Rows included and excluded.
  \item Features included.
  \item Train/test or cross-validation design.
  \item Baseline model.
  \item Candidate model.
  \item Metrics.
  \item Diagnostic plots.
  \item Known limitations.
  \item Use restrictions.
  \item Recommended next validation step.
\end{itemize}

\section{Interpretation Language}

Good model wording:
\begin{itemize}
  \item ``Under the recorded time-based split, the baseline model predicted salinity with lower test-set MAE than the mean-only baseline.''
  \item ``The feature importance summary is treated as a hypothesis-generating screen because correlated predictors and site structure can distort apparent importance.''
  \item ``The classification threshold favors sensitivity, which increases false positives and is appropriate only for screening.''
\end{itemize}

Weak model wording:
\begin{itemize}
  \item ``The model proves which variable causes poor water quality.''
  \item ``High accuracy means the model is ready for field deployment.''
  \item ``This feature is important, so it must be the biological mechanism.''
\end{itemize}
